\chapter*{Methodological Note}
\addcontentsline{toc}{chapter}{Methodological Note}

The dependency rule of this programme is
\[
\boxed{\text{derive temporal structure first; introduce spacetime later}.}
\]
The construction begins from relational information and phase transport. Each new temporal object is admitted only after its upstream dependencies are explicit. The intended order is
\[
\begin{aligned}
\text{information}
&\to\text{directed transition}
\to\text{phase structure}
\to\text{temporal state}\\
&\to\text{NOW/bifurcation}
\to\text{transport/memory}
\to\text{metric calibration}\\
&\to\text{space}
\to\text{Einstein closure}.
\end{aligned}
\]

Mathematical construction, computational evidence, and physical claim status are kept distinct. A candidate is first assigned typed variables, equations, dependencies, and falsification conditions. Algebraic identities are proved where possible; numerical controls are used where analytic closure is not yet available. Only then can a claim move to a stronger registry status.

The ordering parameter used in the early formalism is kept distinct from metric clock time. Metric calibration is therefore a downstream bridge rather than a hidden starting assumption. Likewise, chirality, topology, fractal scaling, memory, and subjective temporal accumulation are treated as derived-property targets with their own observables and tests.

\section*{Prior-art and novelty discipline}

Whenever the construction invokes an established object, the corresponding chapter cites the standard source or a primary historical reference. Shannon entropy, Kullback--Leibler divergence, Pancharatnam/Berry phase, Onsager response, reversible Markov entropy gradients, Noether conservation, Einstein gravity, Bell/CHSH tests, and Hawking/Gibbons--Hawking horizon thermodynamics are therefore treated as external scientific context rather than as IDT discoveries.

A statement is called \emph{IDT-specific} only when it depends on the programme's own typed composition, decoder, activity measure, memory carrier, calibration map, or falsification rule. Recovery of a known physical equation is labelled a recovery benchmark. A physical discriminator requires an independently determined IDT quantity that can be evaluated before the target measurement is inspected.

\section*{Falsification before promotion}

For a proposed physical correspondence, the calibration map and null model must be fixed independently of the target outcome. If the same measured quantity is used both to define the IDT predictor and to confirm it, the exercise is classified as calibration or consistency checking rather than discrimination. Promotion to a physical result requires immutable raw data, uncertainty propagation, declared exclusions, and independent replication where the claim warrants it.
